\section{Omitted Proofs}
\label{app:omitted_proofs}

This appendix collects the proofs omitted from the main body of the paper.
Appendix~\ref{app:ca_local_proof} proves Proposition~\ref{prop:ca_local} of Section~\ref{sec:comparative_advantage}.
Appendix~\ref{app:shortrun_proof} proves Proposition~\ref{prop:shortrun_dp} of Section~\ref{sec:extensions.computation}; the companion long-run result, Proposition~\ref{prop:totalcost_optimization_dp}, is proved in Appendix~\ref{app:jobdesign.longrun}.

\subsection{Proof of Proposition~\ref{prop:ca_local}}
\label{app:ca_local_proof}

Throughout, write $\minTime{j} = \min\left\{\manualTime{j}, \frac{\AItime{j}}{q_j}\right\}$ for the cheapest way of executing step $j$ on its own, whether manually or as a standalone augmented step, and take $q_{k-1}, q_k, q_{k+1} \in (0,1)$.
Because chains are contiguous and none crosses the boundary of the block $\{k-1,\, k,\, k+1\}$, an arrangement of the block is a partition of it into consecutive runs, each run executed either as a single step at cost $\minTime{j}$ or as an AI chain augmented at its last step $r$, at cost $\AItime{r}/\prod_i q_i$ taken over the run's steps.

Before comparing them, note which of these arrangements can have AI execute step $k$ at all.
Executed as a singleton, step $k$ costs $\minTime{k} = \manualTime{k}$, since $\manualTime{k} < \AItime{k}/q_k$ by hypothesis, so augmenting it on its own is strictly dominated by performing it manually and never occurs in an optimum.
AI therefore executes step $k$ only if step $k$ belongs to a chain of length at least two, and since chains are contiguous and confined to the block, that chain is $\{k,\, k+1\}$, $\{k-1,\, k,\, k+1\}$, or $\{k-1,\, k\}$: in the first two, step $k$ is not the chain's last step, so it is never prompted or verified and is automated, while in the third it is the last step, hence augmented.

Pricing a singleton $\{k\}$ at $\manualTime{k}$ accordingly, the four partitions cost
\begin{align*}
V_0 &= \minTime{k-1} + \manualTime{k} + \minTime{k+1}, && \text{(no chain covers step $k$)} \\
V_1 &= \minTime{k-1} + \frac{\AItime{k+1}}{q_k\, q_{k+1}}, && \text{(chain $\{k,\, k+1\}$)} \\
V_2 &= \frac{\AItime{k+1}}{q_{k-1}\, q_k\, q_{k+1}}, && \text{(chain $\{k-1,\, k,\, k+1\}$)} \\
V_3 &= \frac{\AItime{k}}{q_{k-1}\, q_k} + \minTime{k+1}. && \text{(chain $\{k-1,\, k\}$)}
\end{align*}
AI executes step $k$ exactly when $\min\{V_1, V_2, V_3\} < V_0$, ties being resolved in favor of manual execution.

\begin{proof}[Proof of part (i)]
Fix $(\manualTime{k}, \AItime{k}, q_k)$ with $\manualTime{k} < \AItime{k}/q_k$ and fix $q_{k-1}, q_{k+1} \in (0,1)$.
Restrict attention to the two-parameter family of neighbor costs given by $\manualTime{k-1} = \AItime{k-1} = \mu > 0$ and $\manualTime{k+1} = 1$, $\AItime{k+1} = \tau > 0$, so that $\minTime{k-1} = \mu$ and $\minTime{k+1} = \min\{1, \tau/q_{k+1}\}$.
Since $V_0, \dotsc, V_3$ are continuous in all of the neighbor parameters, it suffices to exhibit, for each configuration, values of $(\mu, \tau)$ at which its cost is strictly below the other three; strictness then extends the conclusion to an open neighborhood.

\begin{itemize}
\item \emph{Configuration $\{k,\, k+1\}$, which automates step $k$.}
Set $\tau = q_k\, q_{k+1} \manualTime{k}/2$, so that $V_1 = \mu + \manualTime{k}/2$, and take
\[ 0 < \mu < \min\left\{ \frac{\manualTime{k}}{2}\cdot\frac{1-q_{k-1}}{q_{k-1}},\ \frac{\AItime{k}}{q_{k-1}\, q_k} - \frac{\manualTime{k}}{2} \right\}, \]
a nonempty interval because $q_{k-1} < 1$ and $\AItime{k}/(q_{k-1} q_k) > \AItime{k}/q_k > \manualTime{k}$.
The three comparisons are then
\begin{align*}
V_0 - V_1 &= \frac{\manualTime{k}}{2} + \minTime{k+1}, \\
V_2 - V_1 &= \frac{\manualTime{k}}{2}\cdot\frac{1-q_{k-1}}{q_{k-1}} - \mu, \\
V_3 - V_1 &= \frac{\AItime{k}}{q_{k-1}\, q_k} - \frac{\manualTime{k}}{2} - \mu + \minTime{k+1},
\end{align*}
and all three are positive: the first because $\minTime{k+1} > 0$, and the second and third by the first and second bounds on $\mu$ respectively.

\item \emph{Configuration $\{k-1,\, k,\, k+1\}$, which also automates step $k$.}
Set $\mu = 1$ and take $0 < \tau < q_{k+1} \min\{q_{k-1}\, q_k,\ \AItime{k}\}$, so that the long chain is cheap.
The three comparisons are
\begin{align*}
V_0 - V_2 &= 1 + \manualTime{k} + \minTime{k+1} - \frac{\tau}{q_{k-1}\, q_k\, q_{k+1}}, \\
V_1 - V_2 &= 1 + \frac{\tau}{q_k\, q_{k+1}} - \frac{\tau}{q_{k-1}\, q_k\, q_{k+1}}, \\
V_3 - V_2 &= \frac{\AItime{k}}{q_{k-1}\, q_k} + \minTime{k+1} - \frac{\tau}{q_{k-1}\, q_k\, q_{k+1}}.
\end{align*}
The bound on $\tau$ holds $\tau/(q_{k-1} q_k q_{k+1})$ below both $1$ and $\AItime{k}/(q_{k-1} q_k)$, which leaves all three positive.

\item \emph{Configuration $\{k-1,\, k\}$, which augments step $k$.}
Set $\mu = \AItime{k}/(q_{k-1} q_k)$ and take $\tau > \max\{q_{k+1},\ q_k q_{k+1},\ q_{k-1} q_k q_{k+1} (\mu+1)\}$.
The first of these bounds puts $\tau/q_{k+1}$ above $1$, so that $\minTime{k+1} = 1$ and $V_3 = \mu + 1$, and the comparisons are
\begin{align*}
V_0 - V_3 &= \manualTime{k}, \\
V_1 - V_3 &= \frac{\tau}{q_k\, q_{k+1}} - 1, \\
V_2 - V_3 &= \frac{\tau}{q_{k-1}\, q_k\, q_{k+1}} - (\mu + 1),
\end{align*}
positive because $\manualTime{k} > 0$ and by the second and third bounds on $\tau$ respectively.
\end{itemize}

The margin of the human's advantage on step $k$, $\AItime{k}/q_k - \manualTime{k}$, was unrestricted throughout, which gives the second claim of part~(i).
\end{proof}

\begin{proof}[Proof of part (ii)]
Fix step $k$'s own parameters together with the four cost parameters of its neighbors, and let the two neighbor success probabilities vary over $Q = (0,1) \times (0,1)$.
Write
\[ \mathcal{A} \;=\; \bigl\{ (q_{k-1},\, q_{k+1}) \in Q \; : \; \text{AI executes step } k \bigr\} \]
for the set of neighbor reliabilities at which the cost-minimizing arrangement does not leave step $k$ to the human.
Part~(ii) is precisely the statement that $\mathcal{A}$ is \emph{upward closed} in the coordinatewise order: whenever a pair belongs to $\mathcal{A}$, so does every pair at least as large in both coordinates.
The proof turns that statement into one about the three differences
\[ D_j \;=\; V_j - V_0, \qquad j = 1,\, 2,\, 3, \]
each of which compares one arrangement that has AI execute step $k$ against the arrangement that does not.
Three reductions get us there.

\emph{Reduction 1: from $\mathcal{A}$ to one chain at a time.}
AI executes step $k$ exactly when $\min\{V_1, V_2, V_3\} < V_0$, that is, when $D_j < 0$ for at least one $j$.
Writing $\mathcal{A}_j = \{(q_{k-1},\, q_{k+1}) \in Q : D_j < 0\}$ for the set on which chain $j$ by itself beats manual execution of step $k$, this says
\[ \mathcal{A} \;=\; \mathcal{A}_1 \cup \mathcal{A}_2 \cup \mathcal{A}_3 . \]
A union of upward-closed sets is upward closed, because a pair in the union lies in one of them and every larger pair then lies in that same one.
It is therefore enough to show that each $\mathcal{A}_j$ is upward closed, and the three can be handled separately: we never need to know which chain is optimal, only whether each one, on its own, is cheaper than leaving the step to the human.

\emph{Reduction 2: from both coordinates at once to one at a time.}
Suppose each $\mathcal{A}_j$ is upward closed in $q_{k-1}$ at every fixed value of $q_{k+1}$, and upward closed in $q_{k+1}$ at every fixed value of $q_{k-1}$.
Given $(q_{k-1},\, q_{k+1}) \in \mathcal{A}_j$ and a larger pair $(q'_{k-1},\, q'_{k+1})$, raise the two coordinates in turn:
\[ (q_{k-1},\, q_{k+1}) \in \mathcal{A}_j \;\implies\; (q'_{k-1},\, q_{k+1}) \in \mathcal{A}_j \;\implies\; (q'_{k-1},\, q'_{k+1}) \in \mathcal{A}_j , \]
the first step by closure in the first coordinate holding $q_{k+1}$ fixed, the second by closure in the second coordinate holding $q'_{k-1}$ fixed.
The one-coordinate statements deliver the joint one, provided each holds at every value of the coordinate being held fixed, which is how they are proved below.

\emph{Reduction 3: from upward closure to a property of $D_j$.}
With one coordinate fixed, $D_j$ is a function of a single success probability $q$, and $\mathcal{A}_j$ is the set on which that function is negative.
That set is upward closed as soon as $D_j$, once negative, stays negative as $q$ rises.
It suffices for $D_j$ to be nonincreasing in $q$, which is what we verify in every case but one; in the exceptional case $D_j$ increases, and we show instead that it is negative on the entire range where the increase can occur.

Writing the three differences out,
\begin{align*}
D_1 &= \frac{\AItime{k+1}}{q_k\, q_{k+1}} - \manualTime{k} - \minTime{k+1}, \\
D_2 &= \frac{\AItime{k+1}}{q_{k-1}\, q_k\, q_{k+1}} - \minTime{k-1} - \manualTime{k} - \minTime{k+1}, \\
D_3 &= \frac{\AItime{k}}{q_{k-1}\, q_k} - \minTime{k-1} - \manualTime{k},
\end{align*}
the terms that move with the neighbor reliabilities are the chain costs and the two standalone costs
\[ \minTime{k-1} = \min\left\{ \manualTime{k-1},\, \frac{\AItime{k-1}}{q_{k-1}} \right\}, \qquad \minTime{k+1} = \min\left\{ \manualTime{k+1},\, \frac{\AItime{k+1}}{q_{k+1}} \right\}, \]
of which the first depends on $q_{k-1}$ alone and the second on $q_{k+1}$ alone.

Both coordinates are covered by one observation about functions of a single success probability, which we record first.

\emph{A claim covering both coordinates.}
Let $m,\, A,\, B,\, c > 0$ and define, for $q \in (0,1)$,
\[ D(q) \;=\; \underbrace{\frac{A}{q}}_{\text{a chain through the neighbor}} \;-\; \underbrace{\min\left\{ m,\, \frac{B}{q} \right\}}_{\text{the neighbor on its own}} \;-\; \underbrace{c}_{\text{fixed}} . \]
Then $\{ q \in (0,1) : D(q) < 0 \}$ is upward closed.
To see this, note that the neighbor's standalone cost is the smaller of a manual cost $m$ that does not move with $q$ and an augmented cost $B/q$ that falls as $q$ rises, the two crossing at $q^{*} = B/m$.
Splitting there,
\[
D(q) \;=\;
\begin{cases}
\dfrac{A}{q} - m - c, & 0 < q \leq q^{*}, \\[2.2ex]
\dfrac{A - B}{q} - c, & q^{*} < q < 1,
\end{cases}
\]
the second line obtained by substituting $B/q$ for the minimum.
The two expressions agree at $q = q^{*}$, where $m = B/q^{*}$, so $D$ is continuous; and if $q^{*} \geq 1$ the second branch is empty and only the first arises.
On the first branch $D$ is strictly decreasing, since
\[ \frac{\mathrm{d}}{\mathrm{d}q}\left( \frac{A}{q} - m - c \right) \;=\; -\frac{A}{q^{2}} \;<\; 0 , \]
so raising $q$ cheapens the chain while leaving untouched the manual cost that the chain displaces.
On the second branch the behavior depends on the sign of $A - B$:
\begin{itemize}
\item if $A \geq B$, then $(A-B)/q$ is nonincreasing in $q$, so $D$ is nonincreasing on the branch;
\item if $A < B$, then $D$ increases with $q$, but $(A-B)/q < 0$ and $-c < 0$, so $D(q) < 0$ at \emph{every} point of the branch.
\end{itemize}
Either way, if $D$ is negative at some point of the second branch it is negative everywhere to the right of it.
Now take $q < q'$ in $(0,1)$ with $D(q) < 0$, and consider the three possible positions of the pair relative to $q^{*}$:
\begin{itemize}
\item if $q' \leq q^{*}$, both lie on the first branch and $D(q') < D(q) < 0$;
\item if $q^{*} < q$, both lie on the second branch, which was just handled;
\item if $q \leq q^{*} < q'$, then $D(q^{*}) \leq D(q) < 0$ by the first branch, and since the second-branch expression agrees with $D$ at $q^{*}$ and the argument above applies to it on $[q^{*},\, 1)$, we get $D(q') < 0$.
\end{itemize}
This establishes the claim.
Its content is that raising a neighbor's reliability always cheapens a chain that runs through that neighbor, and cheapens the neighbor's own execution only once that execution is an augmented one, in which case both costs carry the same factor $1/q$ and fall in the same proportion, so the comparison between them cannot reverse.

\emph{Closure in $q_{k-1}$.}
Fix $q_{k+1}$ and let $q = q_{k-1}$ vary.
The difference $D_1$ does not depend on $q$ at all: the two arrangements it compares, the chain $\{k,\, k+1\}$ and manual execution of step $k$, both leave the predecessor to be executed on its own, so the term $\minTime{k-1}$ appears in $V_1$ and $V_0$ alike and cancels.
A constant function is nonincreasing, so $\mathcal{A}_1$ is upward closed in this coordinate.
The other two differences are instances of the claim, with $B = \AItime{k-1}$ and $m = \manualTime{k-1}$ throughout and
\begin{align*}
D_2: \quad A &= \frac{\AItime{k+1}}{q_k\, q_{k+1}}, & c &= \manualTime{k} + \minTime{k+1}, \\
D_3: \quad A &= \frac{\AItime{k}}{q_k}, & c &= \manualTime{k},
\end{align*}
neither of which moves with $q$.
The claim therefore makes $\mathcal{A}_2$ and $\mathcal{A}_3$ upward closed in $q_{k-1}$.
Nothing here ties $A$ to $B$, so both of the claim's second-branch cases can occur: a predecessor whose own augmented cost falls faster than the cost of the chain it would anchor is exactly the case $A < B$, and it is settled not by monotonicity but by the chain's being cheaper throughout that branch to begin with.

\emph{Closure in $q_{k+1}$.}
The successor coordinate is the counterpart of the predecessor one, and the milder half of it, so the claim disposes of it directly.
The difference $D_3$ now plays the role $D_1$ played above and drops out for the mirror-image reason: the chain $\{k-1,\, k\}$ and manual execution of step $k$ both leave the successor to be executed on its own, so $\minTime{k+1}$ cancels between $V_3$ and $V_0$.
The remaining two differences are again instances of the claim, read in $q = q_{k+1}$, with $B = \AItime{k+1}$ and $m = \manualTime{k+1}$ throughout and
\begin{align*}
D_1: \quad A &= \frac{\AItime{k+1}}{q_k}, & c &= \manualTime{k}, \\
D_2: \quad A &= \frac{\AItime{k+1}}{q_{k-1}\, q_k}, & c &= \minTime{k-1} + \manualTime{k} .
\end{align*}
The one difference from the predecessor case is that $A$ and $B$ are no longer free of each other: a chain terminating at the successor pays that same management cost $\AItime{k+1}$, inflated by the failure risk of the steps ahead of it, so $A/B$ equals $1/q_k$ or $1/(q_{k-1} q_k)$ and is at least $1$.
Only the case $A \geq B$ of the claim can arise, and $\mathcal{A}_1$ and $\mathcal{A}_2$ are upward closed in $q_{k+1}$, with the differences nonincreasing outright.

Each $\mathcal{A}_j$ is therefore upward closed in each coordinate separately, hence upward closed in the coordinatewise order by Reduction~2, and $\mathcal{A}$ is upward closed by Reduction~1.
\end{proof}

\subsection{Proof of Proposition~\ref{prop:shortrun_dp}}
\label{app:shortrun_proof}

\begin{proof}
We first note the following recursive formulation of the optimization problem.
For all $k \leq m$, let $C[k]$ denote the minimum cost of completing a hypothetical workflow that only includes steps $1$ through $k$.
Note that, because of the way we defined $C[k]$, task $k$ cannot be automated in any minimum-cost solution that determines $C[k]$; it can only be augmented or performed manually.
Note also that $C[m]$ is the cost of our desired optimal solution.

We now show how to calculate $C[k]$ recursively.
As a base case we have $C[0] = 0$, as the empty set of steps requires no work to complete.
For $k \geq 1$, $C[k]$ is the lesser of
\[ C[k-1] + \manualTime{k} \]
and
\[ \min_{\ell < k} \Biggl\{ C[\ell] + \frac{\AItime{k}}{\prod_{i=\ell+1}^k q_i} \Biggr\}.
\]
In other words, we either complete step $k$ manually (at time $\manualTime{k}$, separately optimizing over steps $1$ through $k-1$) or we augment step $k$, at expected time $\AItime{k}/\prod_{i=\ell+1}^k q_i$.
In the latter case, we then optimize over the length of the AI chain that ends with the newly-augmented step $k$.
This could be a singleton chain with no automation (corresponding to $\ell = k-1$), or a longer chain that automates one or more steps before step $k$ (corresponding to a choice of $\ell < k-1$).
For any such choice of $\ell$, we then separately optimize the production strategy for tasks $1$ through $\ell$.

Using this formulation, given the values $(C[0], \dotsc, C[k-1])$, we can calculate $C[k]$ in time $O(k)$ by considering each potential choice of $\ell$.
Doing so for each $k = 1, 2, \dotsc, m$ yields the value of $C[m]$ (and the corresponding optimal AI strategy) in total time $O(m^2)$.
\end{proof}

\subsection{Proof of Proposition~\ref{prop:fragmentation}}

Recall that Proposition~\ref{prop:fragmentation} relates the cost of the optimal AI deployment strategy for a step sequence to its fragmentation index.
Intuitively, we expect a workflow in which automatable tasks are clustered together will benefit more from AI deployment than one where steps with high and low levels of automation susceptibility are interspersed.
This is due to the potential benefits from AI chains.

First recall the definition of fragmentation index for a given step sequence $\mathcal{S} = (s_1, \dotsc, s_m)$.
As a reminder, we assume for Proposition~\ref{prop:fragmentation} that $\AItime{i} = 1$ for all $s_i$.
Suppose that each step $s_i$ succeeds independently with probability $q_i$; any step that does not succeed is said to fail.
Write $F$ for the set of steps that fail, and $\mathcal{C} = \{C_1, \dotsc, C_k\}$ for the random variable representing the collection of maximal connected components of non-failed steps.
The weight of each $C_j \in \mathcal{C}$ is defined to be $\omega(C_j) = \min\{ 1, \sum_{s_i \in C_j} \manualTime{i} \}$.
That is, each $C_j$ has weight $1$ unless the sum of the manual costs of its steps is less than $1$ (due to our assumption that AI management costs are normalized to $1$).

Recall that for each $s_i$ we write $\minTime{i} = \min\left\{\manualTime{i}, \frac{\AItime{i}}{q_i}\right\}$.
That is, $\minTime{i}$ is the minimum cost to implement step $s_i$ individually (whether manually or through AI augmentation).
Given a realization of $\mathcal{C}$ and $F$, we define the realized fragmentation to be
\[ \sum_{s_i \in F}\minTime{i} + \sum_{C_j \in \mathcal{C}} \omega(C_j).\]
The \emph{fragmentation index} is defined to be the expected value of the realized fragmentation.

We now fix the sequence $\mathcal{S} = (s_1, \dotsc, s_m)$, write $FI$ for the fragmentation index of $\mathcal{S}$, and let $OPT$ be the cost of the optimal (i.e., cost-minimizing) task structure for those steps.

We begin by showing that $FI$ is not much larger than $OPT$.

\begin{lem}
\label{lem:FI.upper.bound}
    $FI \leq \tfrac{5}{4}OPT$.
\end{lem}
\begin{proof}
Fix an arbitrary task sequence $\T = (T_1, \dotsc, T_n)$.
Partition this into two sets of tasks: $\T_{\manualLetter}$ containing all (singleton) tasks that are completed manually, and $\T_{\AIletter}$ containing all other tasks (consisting of AI-augmented singletons and AI chains).
We first claim that
\[ FI \leq \sum_{T_j \in \T_{\AIletter}}\left(1 + \sum_{s_i \in T_j}(1-\alpha^{d_i})(1+\minTime{i})\right) + \sum_{(s_i) \in \T_{\manualLetter}}\manualTime{i}. \]

To see why, first note that for any human-executed task $(s_i) \in \T_{\manualLetter}$, either the task fails in the realization of $F$ (in which case it contributes $\minTime{i} \leq \manualTime{i}$ to $FI$), or it succeeds (in which case it appears in some connected component $C_j$, contributing at most $\manualTime{i}$ to the component's weight $\omega(C_j)$).

Next consider the AI-assisted tasks.
For any $T_j \in \T_{\AIletter}$, consider the set of steps in $T_j$ that fail in any given realization of $F$.
If no steps in $T_j$ fail, then $T_j$ collectively contributes at most $1$ to the realized fragmentation (since all $s_i \in T_j$ lie in the same connected component of non-failed tasks).
Otherwise, each failed task $s_i \in T_j$ contributes $\minTime{i}$ to the realized fragmentation (for itself, recalling that $\minTime{i}$ is the minimal cost of executing $s_i$ on its own) plus an additional value of at most $1$ for the weight of a potential new connected component of non-failed tasks.
As each task $s_i \in T_j$ fails independently with probability $(1-\alpha^{d_i})$, we get that the contribution of tasks in $T_j$ to the fragmentation index is at most $1 + \sum_{s_i \in T_j}(1-\alpha^{d_i})(1+\minTime{i})$ as claimed.

On the other hand, recall that if we take $\T$ to be the cost-minimizing task sequence, then
\[ OPT = \sum_{T_j \in \T_{\AIletter}} \alpha^{-d(T_j)} + \sum_{(s_i) \in \T_{\manualLetter}} \manualTime{i} \]
by definition, where we write $d(T_j) = \sum_{s_i \in T_j}d_i$ for the total difficulty of steps in task $T_j$.

Comparing our expression for $OPT$ with our bound on $FI$, we see that each task in $\T_{\manualLetter}$ contributes the same amount to each, whereas each $T_j \in \T_{\AIletter}$ contributes $1+\sum_{s_i \in T_j}(1-\alpha^{d_i})(1+\minTime{i})$ to the former and $\alpha^{-d(T_j)}$ to the latter.
We will show that, for each $T_j \in \T_{\AIletter}$, $1 + \sum_{s_i \in T_j}(1-\alpha^{d_i})(1+\minTime{i}) \leq \frac{5}{4} \alpha^{-d(T_j)}$, which will complete the proof.

First, since $\minTime{i} \leq \alpha^{-d_i}$ for each $s_i \in \mathcal{S}$, we have $1 + \sum_{s_i \in T_j}(1-\alpha^{d_i})(1+\minTime{i}) \leq 1 + \sum_{s_i \in T_j}(1-\alpha^{d_i})(1+\alpha^{-d_i})$.

Next note that for any $x,y\in [0,1]$, $(1-xy)(1+\tfrac{1}{xy}) \geq (1-x)(1+\tfrac{1}{x}) + (1-y)(1+\tfrac{1}{y})$.
(This can be checked by expanding and taking derivatives.)
Repeatedly applying this fact, we get that $1 + \sum_{s_i \in T_j}(1-\alpha^{d_i})(1+\alpha^{-d_i}) \leq 1 + (1-\alpha^{d(T_j)})(1+\alpha^{-d(T_j)}) = 1 + \alpha^{-d(T_j)} - \alpha^{d(T_j)}$ for any $T_j \in \T_{\AIletter}$.
The ratio between this expression and $\alpha^{-d(T_j)}$ is maximized at $\alpha^{-d(T_j)} = 1/2$, achieving a value of $5/4$ as claimed.
\end{proof}

The bound of $\tfrac{5}{4}$ in Lemma~\ref{lem:FI.upper.bound} is tight, as the following example shows.

\begin{example}
\label{ex:FI.tight}
Consider a three-step workflow in which the first and last steps always succeed, while the second succeeds with probability $1/2$ and is twice as costly to perform manually:
\[ \bigl(\manualTime{i},\, \AItime{i},\, q_i\bigr)_{i=1,2,3} = (1,\,1,\,1),\quad (2,\,1,\,\tfrac12),\quad (1,\,1,\,1). \]
The optimal policy automates all three together for an expected cost of $\AItime{3}/(q_1 q_2 q_3) = 1/(1 \cdot \tfrac12 \cdot 1) = 2$, whereas the fragmentation index is $\tfrac12 \times 1 + \tfrac12 \times (1+2+1) = \tfrac52$.
The ratio $FI/OPT$ is therefore $\tfrac54$.
\end{example}

We next argue that $FI$ is not much smaller than $OPT$.
We begin with an analysis under the assumption that $\manualTime{i} \geq 1$ for all $s_i \in \mathcal{S}$.
That is, for each step $s_i$, completing the step manually is at least as costly as managing an AI tool to complete the step in a single attempt (with guaranteed success).

\begin{lem}
\label{lem:FI.lower.bound.4}
    Suppose $\manualTime{i} \geq 1$ for all $s_i \in \mathcal{S}$.
    Then $FI \geq \tfrac{1}{4}OPT$.
\end{lem}
\begin{proof}

    Consider the following task sequence.
    Beginning with the first step $s_1$, consider the maximal contiguous sequence of steps $T$ whose success probability $\alpha^{d(T)}$ is greater than or equal to $1/2$.
    If that set is empty (i.e., the first step has success probability strictly less than $1/2$) then the first step is set to be completed individually, either manually or augmented, whichever is cheaper.
    In this case, we'll say the resulting singleton task is completed \emph{individually}.
    Otherwise, the sequence $T = (s_1, \dotsc, s_\ell)$ (which could still be a singleton task) is added to the task sequence as an AI chain; we call this a \emph{non-individual} task.
    This procedure is then repeated beginning with the next step (which is $s_2$ in the former case, or $s_{\ell+1}$ in the latter case), until all steps have been added to a task.
    Call the resulting task sequence $\T$.

    \medskip

    Let $ALG$ denote the cost of the task sequence $\T$.
    Note that we must have $ALG \geq OPT$.
    We will argue that $FI \geq ALG/4$, which will prove the claim.

    \medskip

    We first define some notation.
    Write $\T_I$ for the set of individual tasks in $\T$ and $\T_{NI}$ for the set of non-individual tasks.
    Note that $\alpha^{d_i} < 1/2$ for all $(s_i) \in \T_I$, by construction.
    Note also that $\T_{\manualLetter} \subseteq \T_I$ (recalling that $\T_{\manualLetter}$ is the set of all singleton tasks executed manually), and this containment may be strict if $\T_I$ contains singleton tasks that are augmented.
    We emphasize that $\T_{NI}$ may still include singleton tasks, but any singleton task $(s_i) \in \T_{NI}$ must have the property that $\alpha^{d_i} \geq 1/2$.

    If the final task $s_m$ from $\mathcal{S}$ is contained in some $T_j \in \T_{NI}$, we say that $T_j$ is the \emph{terminal} task.
    All other tasks in $\T_{NI}$ are non-terminal tasks.
Note that if the final step is a singleton task in $\T_I$ then no task is terminal.
    For each non-terminal task $T_j \in \T_{NI}$, we'll write $\overline{T}_j$ to denote $T_j$ plus the first step that follows $T_j$.
    For example, if $T_j = (s_i, \dotsc, s_\ell)$, then $\overline{T}_j = (s_i, \dotsc, s_\ell, s_{\ell+1})$.
    The set $\T_{NI}$ then has the property that for each task $T_j \in \T_{NI}$ we have $\alpha^{d(T_j)} \geq 1/2$, and for each non-terminal $T_j \in \T_{NI}$ we have $\alpha^{(d(\overline{T}_j))} < 1/2$.

    \medskip

    We are now ready to relate $FI$ and $ALG$.
    Let $k = |\T_{NI}|$ be the number of non-individual tasks.
    Note first that
    \[ ALG = \sum_{T_j \in \T_{NI}} \alpha^{-d(T_j)} + \sum_{(s_i) \in \T_I} \minTime{i} \leq 2k + \sum_{(s_i) \in \T_I} \minTime{i}. \]
    This is because $\alpha^{d(T_j)} \geq 1/2$ for each $T_j \in \T_{NI}$, and hence $\alpha^{-d(T_j)} \leq 2$.

    \medskip

    Next we bound $FI$.
    Note that the assumption $\manualTime{i} \geq 1$ implies that, in any realization of the connected components of non-failed steps $\mathcal{C} = (C_1, \dotsc, C_\ell)$, $\omega(C_j) = 1$ for all $j$.
    This implies that the realized fragmentation is equal to $1$ plus, for each $s_i \in F$ (i.e., each step $s_i$ that fails), a contribution of $\minTime{i}$ plus an additional $1$ in the event that $i > 1$ and the task immediately preceding $s_i$ did not fail.
    (This additional cost of $1$ accounts for creating a new connected component of non-failed steps ending with $s_{i-1}$.)

    We claim that $FI \geq k/2 + \frac{1}{2}\sum_{(s_i) \in \T_I} \minTime{i}$.
    We will show this by charging  the realized fragmentation to tasks in $\T_{NI}$ and $\T_I$, so that each $T_j \in \T_{NI}$ is charged at least $1$ with probability at least $1/2$, and each $(s_i) \in \T_I$ is charged at least $\minTime{i}$ with probability at least $1/2$.

    To see this, let $E_j$ denote the event that some step in $\overline{T}_j$ fails, for each non-terminal $T_j \in \T_{NI}$.
    Note that $E_j$ occurs with probability at least $1/2$, since $\alpha^{d(\overline{T}_j)} < 1/2$.
    Also, if $E_j$ occurs, then either a step $s_i \in T_j$ fails, or no step in $T_j$ fails but the step immediately following $T_j$ does.
    In the former case, we will charge the $\minTime{i} \geq 1$ contribution of $s_i$'s failure to $T_j$.
    In the latter case, it must be that the step immediately preceding the failed task did not fail; so we will charge the additional contribution of $1$ from $s_i$'s failure (as described above) to $T_j$.
    Additionally, if there is a terminal $T_j$ in $\T_{NI}$, then we charge the baseline $1$ from the calculation of $FI$ to that terminal $T_j$.
    Aggregating over all these cases, we have that at least $1$ is charged to each set $T_j \in \T_{NI}$ with probability at least $1/2$.
    Finally, for each $(s_i) \in \T_I$ that fails (which occurs with probability at least $1/2$, by construction), we charge the $\minTime{i}$ portion of its contribution to task $(s_i)$.
    We conclude that $FI \geq k/2 + \frac{1}{2}\sum_{(s_i) \in \T_I}\minTime{i}$ as claimed.

    But now since $FI \geq k/2 + \frac{1}{2}\sum_{(s_i) \in \T_I}\minTime{i}$ and $ALG \leq 2k + \sum_{(s_i) \in \T_I}\minTime{i}$, we conclude $FI \geq ALG/4$ as claimed.
\end{proof}

The approximation factor in Lemma~\ref{lem:FI.lower.bound.4} cannot be improved below $\tfrac{2}{3}(2+\sqrt{2}) \approx 2.276$, although we do not have a matching lower bound of $4$.

\begin{example}
\label{ex:FI.lower.gap}
Consider a sequence of $m$ identical steps:
\[ \bigl(\manualTime{i},\, \AItime{i},\, q_i\bigr) = \bigl(\sqrt{2},\, 1,\, 1/\sqrt{2}\bigr) \ \text{ for all $i$}. \]
The task sequence described in the proof of Lemma~\ref{lem:FI.lower.bound.4} handles each task separately, for a total cost of $m \sqrt{2}$, so the optimal task sequence performs at least this well.
The fragmentation index is such that each task contributes $\sqrt{2}$ if it fails, plus $1$ if the preceding task did not fail, for a total contribution of $(1-1/\sqrt{2})(\sqrt{2} + 1/\sqrt{2}) = \tfrac{3}{2}(\sqrt{2}-1) \approx 0.6213$.
As $m$ grows large, the ratio between $m\sqrt{2}$ and $1 + m \times 0.6213$ approaches $\tfrac{2}{3}(2+\sqrt{2}) \approx 2.276$.
\end{example}

We can relax the assumption in Lemma~\ref{lem:FI.lower.bound.4} that $\manualTime{i} \geq 1$ for all $s_i$, but at the cost of a worse approximation factor.

\begin{lem}
\label{lem:FI.lower.bound.8}
    For general $\manualTime{i}$ (i.e., allowing $\manualTime{i} < 1$ for some steps $s_i$), $FI \geq \frac{1}{8}OPT$.
\end{lem}
\begin{proof}
    The argument follows the proof of Lemma~\ref{lem:FI.lower.bound.4}, with the following changes.

    First we make a slight change in the definition of the task sequence $\T$ that defines $ALG$.
    If a constructed task $T_j$ has the property that $\sum_{s_i \in T_j} \manualTime{i} < 1$ (which implies it is cheaper to run all steps of $T_j$ manually than as an AI chain that is guaranteed to succeed), we switch to running the tasks of $T_j$ manually in the task sequence.
    In our analysis, we still consider $T_j$ to be a set in $\T_{NI}$; but its cost is taken to be $\sum_{s_i \in T_j} \manualTime{i}$.

    This modification changes our upper bound on the total cost of $ALG$ to the following:
    \[ ALG \leq 2 \sum_{T_j \in \T_{NI}} \min\left\{1, \sum_{s_i \in T_j} \manualTime{i} \right\} + \sum_{ (s_i) \in \T_I} \minTime{i}. \]

    Then, to bound $FI$, we claim that
    \[ FI \geq \frac{1}{4}\sum_{T_j \in \T_{NI}} \min\left\{1, \sum_{s_i \in T_j} \manualTime{i} \right\} + \frac{1}{2}\sum_{(s_i) \in \T_I} \minTime{i} \]
    which would prove the claim.

    To see why this bound on $FI$ holds, we employ a charging argument as before.
    Recall that for each non-terminal $T_j \in \T_{NI}$, the event $E_j$ occurs with probability at least $1/2$.
    When it does, we charge to $T_j$ the contribution $\minTime{i}$ to $FI$ from any $s_i \in F$ that lies in $T_j$.
    Furthermore, when $E_j$ occurs, for each connected component $C_\ell$ that intersects $T_j$ we additionally charge the contribution $\omega(C_\ell)$ from $FI$ to $T_j$.
    Crucially, the contribution from each $C_\ell$ can be charged to at most two different tasks $T_j$ in this way: once for the leftmost $T_j$ that it intersects, and once for the rightmost $T_j$ that it intersects.
    (The reason is that if some $T_j$ is a subset of $C_\ell$, then by definition no step of $T_j$ fails and hence $E_j$ did not occur.
    So this charging can only occur for a task $T_j$ that intersects $C_\ell$ but is not a subset of $C_\ell$, of which there are at most two.)
    Moreover, this charging to $T_j$ adds up to a  total value of at least $\min\left\{1, \sum_{s_i \in T_j} \manualTime{i} \right\}$.
    Since the charging occurs with probability $1/2$ for each $T_j \in \T_{NI}$, and we are at most double-counting each $\omega(C_\ell)$, we conclude that the expected sum of all $\omega(C_\ell)$, plus $\minTime{i}$ for all failed tasks $s_i$ that lie in some $T_j \in \T_{NI}$, is at least $1/4$ of $\sum_{T_j \in \T_{NI}}\min\left\{1, \sum_{s_i \in T_j} \manualTime{i} \right\}$.

    The charging for $(s_i) \in \T_I$ remains unchanged: the value $\minTime{i}$ is charged in the event that step $s_i$ fails, which occurs with probability at least $1/2$.

    Taken together, we obtain our desired $8$ approximation.
\end{proof}

We suspect the factor of $8$ in Lemma~\ref{lem:FI.lower.bound.8} is not tight.
But, as the following example shows, the factor is strictly greater than the factor $4$ from Lemma~\ref{lem:FI.lower.bound.4}, so the assumption that $\manualTime{i} \geq 1$ for all $s_i$ is necessary for Lemma~\ref{lem:FI.lower.bound.4} to hold.

\begin{example}
\label{ex:FI.necessity}
Consider a sequence that alternates between two step types, with $K > 1$ such pairs (so $2K$ steps in total), where $\epsilon$ is vanishingly small:
\[ \bigl(\manualTime{i},\, \AItime{i},\, q_i\bigr) = \bigl(0,\, 1,\, 1/\sqrt{2}-\epsilon\bigr) \ \text{ for odd $i$}, \qquad (1,\,1,\,1) \ \text{ for even $i$}. \]
The task sequence described in the proof of Lemma~\ref{lem:FI.lower.bound.8} groups the steps into pairs of two-step tasks, for a total cost of $K \sqrt{2}$ (ignoring $\epsilon$).
The fragmentation index is $1$ plus $1$ for each task that fails, which is $1 + K(1 - 1/\sqrt{2})$ (again ignoring the impact of $\epsilon$).
As $K$ grows large, the ratio tends to $2(\sqrt{2}+1)$, which is strictly greater than $4$.
\end{example}

\begin{proof}[Proof of Proposition~\ref{prop:fragmentation}]
Combining Lemmas~\ref{lem:FI.upper.bound}, \ref{lem:FI.lower.bound.4}, and~\ref{lem:FI.lower.bound.8} yields Proposition~\ref{prop:fragmentation}: in general $\tfrac{1}{8}OPT \leq FI \leq \tfrac{5}{4}OPT$, and under the additional assumption $\manualTime{i} \geq 1$ for all $i$, $\tfrac{1}{4}OPT \leq FI \leq \tfrac{5}{4}OPT$.
\end{proof}

\subsection{Non-Monotone Gains from AI Improvements: Formal Statement and Proof}
\label{app:nonmonotone}

\begin{table}[!t]
\centering
\caption{Configurations, Costs, and Marginal Gains in Example~\ref{ex:nonmonotone}}
\label{tab:nonmonotone_costs}
\singlespacing
\setlength{\tabcolsep}{10pt}
\begin{tabular}{lccc}
\toprule
Configuration & Time cost ($C$) & $-\mathrm{d}C/\mathrm{d}\alpha$ & Optimal for \\
\midrule
Both steps manual & $\manualTime{1} + \manualTime{2} = 14$ & $0$ & $\alpha \leq 0.50$ \\
\midrule
Step 1 augmented & $\AItime{1}/q_1 + \manualTime{2} = 3.5\alpha^{-11} + 8$ & $38.5/\alpha^{12}$ & never \\
\midrule
Step 2 augmented & $\manualTime{1} + \AItime{2}/q_2 = 6 + 4/\alpha$ & $4/\alpha^{2}$ & $0.50 \leq \alpha \lesssim 0.92$ \\
\midrule
Both steps augmented & $\AItime{1}/q_1 + \AItime{2}/q_2 = 3.5\alpha^{-11} + 4/\alpha$ & $38.5/\alpha^{12} + 4/\alpha^{2}$ & never \\
\midrule
Steps 1--2 chained & $\AItime{2}/(q_1 q_2) = 4\alpha^{-12}$ & $48/\alpha^{13}$ & $\alpha \gtrsim 0.92$ \\
\bottomrule
\end{tabular}
\par\vspace{\fignotesskip}
\begin{minipage}{1\linewidth}
\begin{spacing}{0.2}
{\fontsize{9.5pt}{10pt}\selectfont Notes: The two-step workflow of Example~\ref{ex:nonmonotone}, the AI-hard and AI-easy steps of Example~\ref{ex:ca_predecessor}, in which $\manualTime{1}=6$, $\manualTime{2}=8$, $\AItime{1}=3.5$, $\AItime{2}=4$, and the AI succeeds at step $i$ with probability $q_i = \alpha^{d_i}$ for $d_1 = 11$ and $d_2 = 1$, so that $q_1 = \alpha^{11}$ and $q_2 = \alpha$.
Costs follow from the task time costs of Section~\ref{sec:model.tasks}, and $-\mathrm{d}C/\mathrm{d}\alpha$ is the marginal benefit of improving AI quality.
The last column gives the range of $\alpha$ over which each configuration is cost-minimizing.
Neither configuration in which step~1 is augmented ever minimizes cost.
Augmenting it on its own costs $3.5\alpha^{-11}$, which is prohibitive while AI is unreliable, and by the time AI is reliable enough for that to be affordable, folding step~1 into a chain is cheaper still: an automated step is never prompted or verified on its own, so its management cost is never paid at all.}
\end{spacing}
\end{minipage}
\end{table}


To make the argument of Section~\ref{sec:nonlinear} precise, consider a workflow with any number of steps and chains.
For any AI strategy $\T$, total cost is a fixed manual term plus a polynomial in $1/\alpha$,
\begin{equation}
\label{eq:strategy_cost}
C_\T(\alpha) = K_\T + \sum_{c \in \mathcal{C}(\T)} \AItime{r_c}\,\alpha^{-D_c},
\end{equation}
where $K_\T = \sum_{\text{manual } i} \manualTime{i}$ is the fixed cost of manually-executed steps and the sum runs over the AI chains of $\T$, denoted $\mathcal{C}(\T)$, with each chain $c$ (with augmented endpoint $r_c$) contributing a term of degree equal to its total difficulty $D_c = \sum_{i \in c} d_i \ge 1$.
The cost of $\T$ thus shifts in a continuous and smooth manner as $\alpha$ increases, with monotone marginal gains to technology improvements.

The firm optimizes costs over many possible AI strategies, expressed as the deployment problem \eqref{eq:totalcost_shortterm}.
At each level of AI quality $\alpha$ the firm adopts the cheapest available strategy, so its optimal cost is the lower envelope of the strategy-specific cost curves, $C^*(\alpha) = \min_{\T} C_\T(\alpha)$.
We write $g_\T(\alpha) = -\,\mathrm{d}C_\T/\mathrm{d}\alpha$ for the marginal benefit of improving AI quality under a fixed strategy $\T$, and $g^*(\alpha) = -\,\mathrm{d}C^*/\mathrm{d}\alpha \ge 0$ for the firm's marginal benefit at the optimum; the function $g^*(\alpha)$ equals $g_\T(\alpha)$ on the regime of $\alpha$ values where $\T$ is optimal.
As $\alpha$ rises, the optimal strategy changes at finitely many values of $\alpha$, the reorganization thresholds of Section~\ref{sec:nonlinear}, which separate one regime from the next.
Lemma~\ref{lem:reorg_threshold} shows that although $g^*$ diminishes within each regime it jumps upward at every reorganization threshold, so the returns to improving AI quality are non-monotone.

\begin{lem}[Non-monotone Marginal Returns to AI Quality Improvements]
\label{lem:reorg_threshold}
Let $C^*(\alpha) = \min_{\T} C_\T(\alpha)$ be the firm's optimal cost over all feasible AI strategies and $g^*(\alpha) = -\,\mathrm{d}C^*/\mathrm{d}\alpha$ its marginal benefit of improving AI quality.
Within each regime (an interval of $\alpha$ on which a single AI strategy is optimal), $g^*$ is non-increasing.
Moreover, at every reorganization threshold (where the optimal AI strategy changes), $g^*$ jumps upward.
\end{lem}

\begin{proof}
Each $C_\T$ is a constant plus terms $\AItime{r_c}\,\alpha^{-D_c}$ with $D_c \ge 1$, with derivatives
\[
\begin{aligned}
C_\T'(\alpha) &= -\sum_{c \in \mathcal{C}(\T)} \AItime{r_c}\,D_c\,\alpha^{-D_c-1} \;\le\; 0, \\
C_\T''(\alpha) &= \sum_{c \in \mathcal{C}(\T)} \AItime{r_c}\,D_c(D_c+1)\,\alpha^{-D_c-2} \;\ge\; 0,
\end{aligned}
\]
so each $C_\T$ is non-increasing and convex.

Within a regime where $\T$ is optimal, $g^*(\alpha) = g_\T(\alpha) = -C_\T'(\alpha) \ge 0$, and the convexity of $C_\T$ makes it non-increasing in $\alpha$ (when $\T$ uses no chain $g^*(\alpha) = 0$ but is otherwise strictly decreasing).
The returns to a fixed configuration in a regime thus weakly diminish.
This proves the first claim of the lemma.

Next, let $\alpha_0 \in (0,1)$ be a threshold where the optimum switches from $\T$ (optimal just below $\alpha_0$) to $\T'$ (just above $\alpha_0$).\footnote{We assume this threshold exists. The only scenario where no such threshold exists is the case in which all steps are executed manually for all $\alpha \in (0,1)$. This is the trivial, uninteresting case and is assumed away.}
Define
\[
\phi(\alpha) = C_{\T'}(\alpha) - C_\T(\alpha).
\]
The function $\phi(\alpha)$ is the difference of two finite sums of powers of $\alpha$ and hence continuous and differentiable on $(0,1)$. 
The two costs are equal at $\alpha_0$, with $C_{\T'}$ the cheaper just above and $C_\T$ the cheaper just below.
Therefore, by definition $\phi$ satisfies:
\[
\phi(\alpha_0^{-}) \ge 0,
\qquad
\phi(\alpha_0) = 0,
\qquad
\phi(\alpha_0^{+}) \le 0.
\]
A differentiable function with such a sign change has nonpositive derivative in the local neighborhood of $\alpha_0$.
So $\phi'(\alpha_0) \le 0$, or equivalently, $C_{\T'}'(\alpha_0) \le C_\T'(\alpha_0)$.
By the definition $g_\T(\alpha) = -\,\mathrm{d}C_\T/\mathrm{d}\alpha$, this is
\[
g_\T(\alpha_0) \;=\; -C_\T'(\alpha_0)
\;\le\;
-C_{\T'}'(\alpha_0) \;=\; g_{\T'}(\alpha_0).
\]
Since $g^*(\alpha)$ coincides with $g_\T(\alpha)$ just below $\alpha_0$ and with $g_{\T'}(\alpha)$ just above, the firm's marginal benefit function $g^*(\alpha)$ jumps upward as $\alpha$ crosses $\alpha_0$.
\end{proof}
